%% independent_operators.tex
%% Operator dependence and the rank of Gamma. A note for the maintainers of PyTreeNet and the
%% authors of the symbolic Gaussian elimination paper. Every number is measured by
%% experiments/independent_operators/independent_operators_probe.py and saved beside it
%% in probe_output.txt.
%%
%% Build (latexmk needs Perl, which is not installed here; quantikz needs the extra pass):
%%   pdflatex -synctex=1 -interaction=nonstopmode -file-line-error independent_operators  (three times)

\documentclass[11pt,a4paper]{article}

\usepackage[T1]{fontenc}
\usepackage{lmodern}
\usepackage{microtype}
\usepackage[margin=2.4cm]{geometry}
\usepackage{amsmath,amssymb,mathtools}
\usepackage{booktabs}
\usepackage{tabularx}
\usepackage{array}
\usepackage{tikz}
\usetikzlibrary{arrows.meta,positioning}
\usepackage{quantikz}
\usepackage[font=small,labelfont=bf,skip=6pt]{caption}
\captionsetup[table]{position=top}
\usepackage[section]{placeins}
\usepackage[hidelinks]{hyperref}
\usepackage{cleveref}

\setlength{\parindent}{0pt}
\setlength{\parskip}{0.45em}
\setlength{\emergencystretch}{3em}
\renewcommand{\arraystretch}{1.15}
\newcolumntype{L}{>{\raggedright\arraybackslash}X}

\newcommand{\code}[1]{\texttt{#1}}
\newcommand{\file}[1]{\texttt{\nolinkurl{#1}}}
\newcommand{\vectorise}{\operatorname{vec}}
\newcommand{\rank}{\operatorname{rank}}
\definecolor{bondred}{RGB}{180,40,40}

\begin{document}

\begin{center}
{\LARGE\bfseries Operator dependence and the rank of $\Gamma$\par - AI-Generated note}
\vspace{0.5em}
\end{center}

%% ------------------------------------------------------------------------------------------------
\section*{Introduction}

The quantity your algorithm drives down on a cut is $\rank(\Gamma)$, with $\Gamma$ the matrix of
your Eq.~(4). The sets $U$ and $V$ behind its rows and columns are formed by removing duplicated operator
chains, so their elements are distinct \emph{as symbols}. The identity $\rank(\Gamma) = $ minimal
bond needs more than that. It needs the chains to be linearly independent \emph{as operators}, and
that is a property of the matrices the symbols name, which the elimination does not hold. Where it
fails, the bond is above the operator-Schmidt rank of the operator and no elimination on $\Gamma$
reaches it. A Hamiltonian meets this whenever its local vocabulary is dependent. An XY chain whose
hopping is written with ladder operators and whose anisotropy is written with Cartesian components
comes out at bond $6$ on every cut where the operator-Schmidt rank is $4$. The smallest instance of
the same failure is a controlled gate whose target carries a multiple of the identity.

This note states the condition exactly, gives the counterexample, and adds a pre-processing step
that establishes the condition over the rationals before $\Gamma$ is built. The step is exact and
leaves the operator unchanged. It returns a Hamiltonian whose operators are already independent
untouched, which includes every sum of Pauli strings and so every Hamiltonian of the experiments in
your Sec.~IV~A.

%% ------------------------------------------------------------------------------------------------
\section{When \texorpdfstring{$\rank(\Gamma)$}{rank(Gamma)} is the minimal bond}\label{sec:gamma}

Fix one cut. With the notation of your Sec.~III~A, writing $n = |U|$ and $m = |V|$,
%
\begin{equation}\label{eq:sum}
  \hat{O} \;=\; \sum_{i=1}^{n}\sum_{j=1}^{m} \Gamma_{ij}\, \hat{U}_i \otimes \hat{V}_j .
\end{equation}
%
Let $D_U$ and $D_V$ be the total physical dimensions of the two sides of the cut, let
$\vectorise$ stack the entries of an operator in a fixed order, and collect the chains in
%
\begin{equation}\label{eq:lr}
  L = \begin{bmatrix} \vectorise \hat{U}_1 & \cdots & \vectorise \hat{U}_n \end{bmatrix}
      \in \mathbb{C}^{D_U^2 \times n},
  \qquad
  R = \begin{bmatrix} \vectorise \hat{V}_1 & \cdots & \vectorise \hat{V}_m \end{bmatrix}
      \in \mathbb{C}^{D_V^2 \times m}.
\end{equation}
%
Reshaping $\hat{O}$ across the cut into the matrix $M$ whose row index is the pair of physical
indices on the left and whose column index is the pair on the right, \cref{eq:sum} reads
%
\begin{equation}\label{eq:reshape}
  M \;=\; L\,\Gamma\,R^{\mathsf{T}}.
\end{equation}
%
The smallest bond this cut admits, in any MPO or TTNO that represents $\hat{O}$, is $\rank(M)$.
Together with \cref{eq:reshape} this gives
%
\begin{equation}\label{eq:bound}
  \rank(M) \;\le\; \rank(\Gamma),
\end{equation}
%
with equality when $L$ and $R$ have full column rank, because multiplication by a matrix with
independent columns on either side leaves the rank unchanged. Full column rank of $L$ is exactly the
statement that $\hat{U}_1,\dots,\hat{U}_n$ are linearly independent operators. The Frobenius rank
inequality bounds the gap by the two dependence defects,
%
\begin{equation}\label{eq:frobenius}
  \rank(M) \;\ge\; \rank(\Gamma) - \bigl(n - \rank L\bigr) - \bigl(m - \rank R\bigr).
\end{equation}
%
The condition is decided one node at a time. Suppose every node carries a linearly independent set
of operators. The chains of $U$ are then distinct tensor products of elements of those sets, hence
distinct elements of a basis of the tensor product of their spans, hence linearly independent, so
$L$ has full column rank. The same holds for $R$. Conversely, if the chains of $U$ are dependent,
then some node in the left part carries a dependent set, since otherwise the previous sentence
applies. So
%
\begin{equation}\label{eq:criterion}
  \text{operators independent on every node}
  \;\Longleftrightarrow\;
  \rank(\Gamma) = \rank(M) \text{ on every cut},
\end{equation}
%
and the right side holds for every cut at once or fails only at cuts separated by a node that
carries a dependence. The argument uses only that a cut splits the nodes into two sets, so it covers
a tree exactly as it covers a chain. One consequence is a test that costs nothing. A node whose $p$
distinct symbols are matrices of dimension $d$ carries a dependence as soon as $p > d^2$. The
converse does not hold, since three operators of dimension two can already be dependent, so $p \le
d^2$ screens nothing and a Hamiltonian is immune only when the operators it draws on are themselves
independent.

Your paper states a condition of this shape for the bipartite method of Ren et al.\ \cite{Ren2020},
on the prefactors rather than on
the operators, \emph{``This algorithm can only find the optimal bond dimension if all prefactors
$\gamma_i$ of Eq.~(3) are different and independent''} (p.~4), and for SGE it reports
\emph{``We found that our SGE-based algorithm achieves optimal bond dimensions and outperforms the
original purely bipartite-graph-based algorithm of [11] for a diverse set of Hamiltonians. A
mathematical proof of optimality is still outstanding''} (p.~15). The condition in
\cref{eq:criterion} is on the operators and is independent of the prefactors. It is not established
by the elimination, which stores the entries of $\Gamma$ as pairs $(c, s_i)$ and prohibits the sum of
two different symbols by your Eq.~(8), so combining two dependent operators into one is outside the moves
available to it.

%% ------------------------------------------------------------------------------------------------
\section{A spin chain over a dependent vocabulary}\label{sec:chain}

Take the chain
%
\begin{equation}\label{eq:xy}
  \hat{H} \;=\; \sum_{i} J \bigl( \hat{S}^{+}_i \hat{S}^{-}_{i+1}
                                + \hat{S}^{-}_i \hat{S}^{+}_{i+1} \bigr)
            \;+\; \sum_{i} g \bigl( \hat{X}_i \hat{X}_{i+1} + \hat{Y}_i \hat{Y}_{i+1} \bigr)
            \;+\; \sum_{i} h\, \hat{Z}_i ,
\end{equation}
%
with $J = 1$, $g = 0.4$ and $h = 0.3$. The hopping is written with ladder operators and the
anisotropy with Cartesian components, which is what happens when the two usual forms of the same
chain are both present in one Hamiltonian. An interior node then carries the six symbols
$\hat{S}^{+}$, $\hat{S}^{-}$, $\hat{X}$, $\hat{Y}$, $\hat{Z}$, $\hat{I}$, and
$\hat{S}^{\pm} = (\hat{X} \pm i\hat{Y})/2$, so four of them span two dimensions. The counting test
of \cref{sec:gamma} also flags the node, $p = 6 > d^2 = 4$.

On six sites the construction gives bond $6$ on every cut, where the operator-Schmidt rank of
$\hat{H}$ is $4$. So \cref{eq:bound} is strict, by the dependence defect of the node vocabulary.
Nothing in the input marks the case. Every string is distinct, the prefactors are distinct, and the
chain sets $U$ and $V$ are non-redundant as symbols. Rewriting the four dependent symbols over two
of them, as in \cref{sec:rewrite}, brings every cut to $4$ and takes the number of strings from $26$
to $16$.

%% ------------------------------------------------------------------------------------------------
\section{A controlled gate, the minimal instance}\label{sec:phase}

The same failure at the smallest size it occurs. A controlled gate enters as the three strings of
$\hat{C}_b\hat{U} = \hat{I} + \hat{\Pi}_b \otimes (\hat{U} - \hat{I})$, namely
$\hat{\Pi}_b \otimes \hat{U}$ with prefactor $+1$, $\hat{\Pi}_b \otimes \hat{I}$ with $-1$, and
$\hat{I} \otimes \hat{I}$ with $+1$, where $\hat{\Pi}_b$ is the projector onto the control pattern.
Take a target that carries a phase on both of its levels, $\hat{U} = \lambda \hat{I}$. A controlled
phase gate in the usual sense, $\hat{U} = \mathrm{diag}(1, e^{i\theta})$, is not this case, since
that matrix is independent of $\hat{I}$ and the gate already comes out at its rank.

\begin{figure}[htbp]
\centering
\begin{quantikz}[column sep=10pt, row sep=10pt]
  \lstick{$q_1$} & \ctrl{1} & \qw \\
  \lstick{$q_2$} & \ctrl{2} & \qw \\
  \lstick{$q_3$} & \qw      & \qw \\
  \lstick{$q_4$} & \gate{\lambda} & \qw
\end{quantikz}
\hspace{2.2cm}
\begin{quantikz}[column sep=10pt, row sep=10pt]
  \lstick{$q_1$} & \ctrl{3} & \qw \\
  \lstick{$q_2$} & \qw      & \qw \\
  \lstick{$q_3$} & \qw      & \qw \\
  \lstick{$q_4$} & \ctrl{1} & \qw \\
  \lstick{$q_5$} & \gate{\lambda} & \qw
\end{quantikz}
\caption{Two controlled phases. On the left two controls next to the target, on the right a gap
between the controls. The target carries $\lambda \hat{I}$, which is dependent on the $\hat{I}$ the
other two strings carry there.}
\label{fig:circuits}
\end{figure}

Take the gate on the left of \cref{fig:circuits}, four qubits with controls on $q_1$ and $q_2$ and
the target on $q_4$, and the cut between $q_2$ and $q_3$. The three strings give
$U = \{\hat{\Pi}_b,\ \hat{I}\otimes\hat{I}\}$ and
$V = \{\hat{I}\otimes\hat{U},\ \hat{I}\otimes\hat{I}\}$, both non-redundant as symbols, and
%
\begin{equation}\label{eq:gammaphase}
  \Gamma = \begin{pmatrix} 1 & -1 \\ 0 & 1 \end{pmatrix},
  \qquad \rank(\Gamma) = 2 .
\end{equation}
%
No row or column operation on \cref{eq:gammaphase} lowers that. The two elements of $V$ are
$\hat{I}\otimes\hat{U} = \lambda\, \hat{I}\otimes\hat{I}$, so $\rank R = 1$, and \cref{eq:reshape}
gives $\rank(M) = 1$. The gate is supported on its controls alone,
$\hat{C}_b\hat{U} = \hat{I} + (\lambda - 1)\, \hat{\Pi}_b \otimes \hat{I}$, and one bond state
carries it. The built bonds are $[2,2,2]$ where the operator-Schmidt rank is $[2,1,1]$.

%% ------------------------------------------------------------------------------------------------
\section{Rewriting over an independent set}\label{sec:rewrite}

On node $k$ let the distinct symbols be $s_1,\dots,s_p$ with matrices $A_1,\dots,A_p$. Choose
$i_1 < \dots < i_q$ so that $A_{i_1},\dots,A_{i_q}$ is a basis of the span of all of them, and write
each operator in that basis,
%
\begin{equation}\label{eq:coords}
  A_t \;=\; \sum_{b=1}^{q} c^{(k)}_{tb}\, A_{i_b},
  \qquad t = 1,\dots,p .
\end{equation}
%
Substituting \cref{eq:coords} into a term of your Eq.~(3) and expanding the product over the nodes,
%
\begin{equation}\label{eq:expand}
  \gamma \bigotimes_{k} A_{t_k}
  \;=\; \sum_{b_1,\dots,b_N}\;
        \gamma \Bigl(\prod_{k} c^{(k)}_{t_k b_k}\Bigr) \bigotimes_{k} A_{i_{b_k}},
\end{equation}
%
and strings that have become equal are collected into one with the sum of their prefactors. The
operator is unchanged, every node now carries only the operators of its own basis, and these are
independent by construction, so \cref{eq:criterion} holds and $\rank(\Gamma)$ is the minimal bond on
every cut. Nodes whose operators are already independent have $q = p$ and $c^{(k)}_{tb} = \delta_{tb}$
in \cref{eq:coords}, so they contribute one summand to \cref{eq:expand} and nothing moves. A
Hamiltonian independent on every node is returned as the object it came in as.

\textbf{Exactness.} A prefactor is carried as $a + ib$ with $a$ and $b$ rational. A float converts to
a rational without rounding, being a dyadic rational already, so no conversion loses a digit. The
elimination that produces \cref{eq:coords} takes the first nonzero entry of a column as its pivot,
with no comparison of magnitudes, which is sound because the arithmetic is exact. No rounding error
can therefore report a dependence that is not there or miss one that is, and the bond is a property
of the operator rather than of the floating point format.

\textbf{The symbol set.} A prefactor that stays rational after \cref{eq:expand} is returned on the
symbol it came in with. Only one that turns complex is given a symbol of its own, and equal values
share it. The set $S$ therefore grows, before the elimination runs, by at most the number of distinct
complex prefactors the rewrite creates, and the restriction of your Eq.~(8) during the elimination is
untouched.

\textbf{Term count.} \Cref{eq:expand} is a product over the nodes, so the rewrite can produce more
strings than it consumed, and collection can leave fewer. Both directions occur, and both are
measured below.

%% ------------------------------------------------------------------------------------------------
\section{When an exact dependence must be refused}\label{sec:refuse}

Over the rationals a set of operators is dependent or it is not, with no tolerance in between. That
is the property that makes the bond exact, and it is also a trap. Operators that are dependent only
to within rounding are exactly independent, and operators that are exactly dependent on a basis that
is nearly rank deficient have coordinates in \cref{eq:coords} of size $\lVert A_t\rVert / \sigma_q$
with $\sigma_q$ the smallest singular value of the basis. Evaluating $\sum_b c_{tb} A_{i_b}$ in
double precision then loses the operator to cancellation.

Five diagonal operators of dimension four, four of which agree to $10^{-9}$, have exact rank four
and numerical rank three. The singular values of the first four are $4.15$, $5.3\cdot 10^{-9}$,
$3.2\cdot 10^{-16}$ and $5.3\cdot 10^{-17}$, the exact coordinates of the fifth on them reach
$4.8\cdot 10^{16}$, and evaluating \cref{eq:coords} in double precision misses the operator by
$11.4$ times its largest entry.

A dependence found exactly is therefore used only when it also holds in double precision. It must
rebuild the operator to a relative $10^{-10}$, and the kept operators must have numerical rank $q$ at
the same tolerance. A node that fails either test is left as it is, which costs a bond of the TTNO
and never the operator it represents. This is what keeps a generator whose entries are floats, such
as a Lindbladian, correct under the pass.

%% ------------------------------------------------------------------------------------------------
\section{What changes}\label{sec:changes}

Every number below is produced by \file{experiments/independent_operators/independent_operators_probe.py} and saved in
\file{probe_output.txt} beside it. Each case is built twice, once with the rewrite and once without,
the two dense operators are compared, and the rank column is the operator-Schmidt rank from a
singular value decomposition of that operator. Bonds are listed left to right along the chain.

\begin{table}[htbp]
\caption{Cases in which the bond falls. The controlled phases carry $\lambda = e^{0.7i}$ and the
sweep behind the first three rows runs one to three controls and two to eight sites, nineteen cases,
every one of which lands on the rank.}
\label{tab:improved}
\centering
\small
\begin{tabular}{llll}
\toprule
case & without the rewrite & with it & rank \\
\midrule
one control, 8 sites           & $2,2,2,2,2,2,2$ & $1,1,1,1,1,1,1$ & $1,1,1,1,1,1,1$ \\
two controls, 8 sites          & $2,2,2,2,2,2,2$ & $2,1,1,1,1,1,1$ & $2,1,1,1,1,1,1$ \\
three controls, 8 sites        & $2,2,2,2,2,2,2$ & $2,2,1,1,1,1,1$ & $2,2,1,1,1,1,1$ \\
controls on 1 and 4, 8 sites   & $2,2,2,2,2,2,2$ & $2,2,2,1,1,1,1$ & $2,2,2,1,1,1,1$ \\
controlled identity, 6 sites   & $2,2,2,2,2$     & $1,1,1,1,1$     & $1,1,1,1,1$ \\
$\hat{I}\hat{X} - \hat{P}_0\hat{X} - \hat{P}_1\hat{X} + \hat{P}_0\hat{Z}$
                               & $2$             & $1$             & $1$ \\
$\hat{n}\hat{n}$ coupling with a $\hat{Z}$ field, 6 sites
                               & $3,3,3,3,3$     & $2,3,3,3,2$     & $2,3,3,3,2$ \\
XY chain, ladder and Cartesian, 6 sites
                               & $6,6,6,6,6$     & $4,4,4,4,4$     & $4,4,4,4,4$ \\
Lindbladian, two jumps, 2 qubits & $5,4,5$       & $4,4,4$         & $4,4,4$ \\
Lindbladian, two jumps, 3 qubits & $5,4,6,4,5$   & $4,4,6,4,4$     & $4,4,5,4,4$ \\
\bottomrule
\end{tabular}
\end{table}

The first four rows are the gate of \cref{sec:phase}, and the bond is $1$ on every cut outside the
span of the controls, which is where the gate acts.

The rows that matter in practice are the next three, where the dependence comes from the vocabulary
a Hamiltonian is written in rather than from a gate. The seventh row is a chain with coupling
$\hat{n}\otimes\hat{n}$ and field $\hat{Z}$, where $\hat{n} = (\hat{I} - \hat{Z})/2$, so the three
symbols $\hat{n}$, $\hat{Z}$ and $\hat{I}$ on a node span two dimensions. The eighth is an XY chain
whose hopping is written with $\hat{S}^{+}$ and $\hat{S}^{-}$ and whose anisotropy is written with
$\hat{X}$ and $\hat{Y}$, four symbols on a node spanning two dimensions, and its string count falls
from $26$ to $16$ as well. Writing one chain in two vocabularies is ordinary, and nothing warns the
author that it costs half again on every bond. The last two rows are the symbolic Lindbladian of
\code{generate\_lindbladian} with two jump operators per qubit, on a chain whose nodes are the ket
and the bra copy of each qubit.

\begin{table}[htbp]
\caption{Cases that do not move, which is the argument that the rewrite is safe to run by default.}
\label{tab:unmoved}
\centering
\begin{tabularx}{\textwidth}{>{\raggedright\arraybackslash}p{0.40\textwidth}L}
\toprule
case & result \\
\midrule
randomised Hamiltonians over $\{\hat{X},\hat{Y},\hat{Z},\hat{I}\}$ on the eight-site tree of
your Sec.~IV~A, 30 terms, 20 seeds
  & the same object is returned 20 times out of 20, so the TTNO is identical by construction \\
\code{MultiControlledGate} with $\hat{X}$, $\hat{Y}$, $\hat{Z}$ and $\hat{H}$ on four qubits
  & bonds $2,2,2$ either way, the same object is returned \\
a Lindbladian in the Pauli basis with float entries, three and four qubits
  & the TTNO matches the dense generator to $2.2\cdot 10^{-15}$ and $3.6\cdot 10^{-15}$, the node
    being refused by the test of \cref{sec:refuse} \\
controlled gates on four qubits whose target is $\mathrm{diag}(1, e^{i\theta})$, $\hat{X}$ or
$\hat{H}$
  & bonds $2,2,2$ either way, equal to the rank, and the terms are returned unchanged, those
    targets being independent of $\hat{I}$ \\
one term, and a sum in which every prefactor cancels
  & the pass is the identity and does not return an empty Hamiltonian \\
\bottomrule
\end{tabularx}
\end{table}

The first row is general rather than fortunate. Those Hamiltonians draw their operators from four
matrices of dimension two that form a basis, so every node carries an independent set and
\cref{eq:criterion} already holds. No Hamiltonian built from a pool of independent operators, and in
particular no sum of Pauli strings, can be affected by the rewrite in either direction, and no second
diagram is built for one.

\textbf{What the rewrite does not guarantee.} \Cref{eq:criterion} makes $\rank(\Gamma)$ equal to the
minimal bond. It does not make the elimination attain $\rank(\Gamma)$, and the rewrite changes the
presentation the elimination runs on, so the rewritten terms can come out worse. On the three-qubit
chain of the last row with one jump operator per qubit instead of two, the terms as given reach $5$
on the middle cut, which is the rank, while the terms rewritten reach $6$. Lines 20 and 21 of your
Algorithm~1 keep the better of two results in the same situation, and the pass does the same. Both
diagrams are built whenever the rewrite changes the terms, and the one with the smaller total bond
dimension is kept, so no case measured here comes out above the bond the terms as given would have
produced. What remains is that neither presentation need attain $\rank(\Gamma)$. The middle cut of
the last row of \cref{tab:improved} stands at $6$ against a rank of $5$.

\textbf{Cost.} The rewrite is exact rational arithmetic on $p$ vectors of length $d^2$ per node, and
it ran in under $0.04$\,s on every case above. A Hamiltonian it leaves alone costs one diagram, as
before, and one it changes costs two, which came to between $0.005$\,s and $0.11$\,s in total for
every case in \cref{tab:improved}. The string count can fall or rise. A controlled gate
goes from three strings to two, the chain of the seventh row of \cref{tab:improved} from $17$ to
$12$, and the Lindbladian of the last row from $34$ to $38$, since \cref{eq:expand} is a product over
the nodes. The pass is reached through \code{combine\_dependent\_operators}, which defaults to
\code{True}, and passing \code{False} restores the previous path exactly.


%% ------------------------------------------------------------------------------------------------
{\small
\begin{thebibliography}{9}
\setlength{\itemsep}{1pt}
\bibitem{Cakir2025}
H.~\c{C}ak{\i}r, R.~M. Milbradt and C.~B. Mendl, \emph{Optimal Symbolic Construction of Matrix
Product Operators and Tree Tensor Network Operators}, arXiv:2502.18630.

\bibitem{Ren2020}
J.~Ren, W.~Li, T.~Jiang and Z.~Shuai, \emph{A general automatic method for optimal construction of
matrix product operators using bipartite graph theory}, J.~Chem.\ Phys.\ \textbf{153}, 084118 (2020).
\end{thebibliography}
}

\end{document}
